\documentclass[12pt]{article}

% ============================================================
% IDIOMA Y CODIFICACIÓN
% ============================================================
\usepackage[utf8]{inputenc}
\usepackage[T1]{fontenc}
\usepackage[spanish, es-noshorthands]{babel}

% ============================================================
% COLOR
% ============================================================
\usepackage[dvipsnames]{xcolor}

% ============================================================
% MÁRGENES Y GEOMETRÍA
% ============================================================
\usepackage[a4paper, margin=2.5cm]{geometry}

% ============================================================
% ESPACIADO MEJORADO
% ============================================================
\usepackage{setspace}
\setstretch{1.15}

\setlength{\parskip}{0.5em}
\setlength{\parindent}{0pt}

% ============================================================
% FORMATO DE TÍTULOS
% ============================================================
\usepackage{titlesec}

\titleformat{\section}
  {\Large\bfseries}
  {\thesection.}
  {0.5em}
  {}

\titlespacing*{\section}{0pt}{1.5em}{0.8em}

% ============================================================
% MATEMÁTICAS
% ============================================================
\usepackage{amsmath, amssymb, amsfonts}

% ============================================================
% CÓDIGO FUENTE
% ============================================================
\usepackage{listings}

% ============================================================
% GRÁFICOS
% ============================================================
\usepackage{graphicx}
\usepackage{float}
\usepackage{tikz}
\usetikzlibrary{arrows.meta}
\usepackage{pgfplots}
\pgfplotsset{compat=1.18}
\usepgfplotslibrary{fillbetween}

% ============================================================
% TABLAS
% ============================================================
\usepackage{booktabs}
\usepackage{array}
\usepackage{multirow}

\definecolor{codebg}{RGB}{248, 248, 248}
\definecolor{codeframe}{RGB}{200, 200, 200}
\definecolor{codecomment}{RGB}{106, 153, 85}
\definecolor{codekeyword}{RGB}{86, 156, 214}
\definecolor{codestring}{RGB}{206, 145, 120}
\definecolor{codenumber}{RGB}{181, 206, 168}

% ============================================================
% LISTAS Y NUMERACIÓN
% ============================================================
\usepackage{enumitem}

% ============================================================
% ENLACES E HIPERTEXTO
% ============================================================
\usepackage{hyperref}
\hypersetup{
    colorlinks  = true,
    linkcolor   = black,
    urlcolor    = blue!60!black,
    citecolor   = black,
    pdftitle    = {Trabajo Práctico N°2 - Parte 2},
    pdfauthor   = {Facundo Luna},
}

% ============================================================
% ENCABEZADO Y PIE DE PÁGINA
% ============================================================
\usepackage{fancyhdr}
\pagestyle{fancy}
\fancyhf{}
\renewcommand{\headrulewidth}{0.4pt}
\renewcommand{\footrulewidth}{0.4pt}
\fancyhead[L]{\small Cálculo Diferencial e Integral II}
\fancyhead[R]{\small Trabajo Práctico N°2 - Parte 2}
\fancyfoot[L]{\small Facundo Luna}
\fancyfoot[R]{\small \thepage}

% ============================================================
% DATOS DEL TRABAJO
% ============================================================

\newcommand{\universidad}{Universidad Católica de Salta}
\newcommand{\carrera}{Licenciatura en Ciencia de Datos}
\newcommand{\materia}{Cálculo Diferencial e Integral II}
\newcommand{\trabajo}{Trabajo Práctico N°2 - Parte 2}
\newcommand{\profesor}{Lic. Mauro Vaca}
\newcommand{\alumno}{Facundo Luna}
\newcommand{\fecha}{\today}

\newcommand{\logoincluido}{%
    \includegraphics[width=3.5cm]{../../../../assets/ucasal-logo.png}\\[0.5cm]%
}
\newcommand{\logoportada}{\logoincluido}

% ============================================================
% COMANDO SOLUCIÓN
% ============================================================
\newcommand{\solucion}{%
    \vspace{0.8em}%
    \textbf{Solución:}%
}

% ============================================================
% COMANDO VOLVER AL ÍNDICE
% ============================================================
\newcommand{\volverindice}{%
    \vspace{1em}%
    \begin{flushright}%
        \hyperlink{indice}{\small $\uparrow$ Volver al índice}%
    \end{flushright}%
}

% ============================================================

\begin{document}

% ======================== PORTADA ===========================

\begin{titlepage}
    \centering

    \logoportada

    \rule{\linewidth}{1.5pt}\\[0.4cm]

    {\large\scshape \universidad}\\[0.2cm]
    {\large \carrera}\\[0.1cm]

    \rule{\linewidth}{0.5pt}\\[1.5cm]

    {\Huge\bfseries \trabajo}\\[0.8cm]
    {\LARGE \materia}\\[2cm]

    \begin{tabular}{rl}
        \textbf{Alumno:}   & \alumno   \\[0.3cm]
        \textbf{Profesor:} & \profesor \\[0.3cm]
        \textbf{Fecha:}    & \fecha    \\
    \end{tabular}

    \vfill

    \rule{\linewidth}{1.5pt}

\end{titlepage}

% ======================== ÍNDICE ============================

\newpage
\tableofcontents
\hypertarget{indice}{}
\setcounter{tocdepth}{2}

% La parte 2 continúa la numeración de la parte 1 (ejercicios 1 a 8).
\setcounter{section}{8}

% ======================== EJERCICIO 9 ========================

\newpage

\section{Propiedades de la integral definida}

Suponiendo conocidos:
\[
\int_{-4}^{5}f(x)\,dx=10,
\qquad
\int_{0}^{3}f(x)\,dx=7,
\qquad
\int_{3}^{5}f(x)\,dx=-4,
\]
\[
\int_{-4}^{-2}f(x)\,dx=-5,
\qquad
\int_{-4}^{-2}g(x)\,dx=-2,
\qquad
\int_{-4}^{0}g(x)\,dx=0.
\]

Utilice propiedades de la integral definida para calcular:

\begin{enumerate}[label=\alph*.-]

    \item
    \[
    \int_{3}^{5}4f(x)\,dx
    \]

    \solucion

    Una constante multiplicativa puede salir fuera de la integral:
    \[
    \int_{3}^{5}4f(x)\,dx=4\int_{3}^{5}f(x)\,dx=4\cdot(-4).
    \]

    \[
    \boxed{
    \int_{3}^{5}4f(x)\,dx=-16
    }
    \]


    \item
    \[
    \int_{-2}^{0}g(x)\,dx+\int_{-2}^{-4}g(x)\,dx
    \]

    \solucion

    Por la aditividad respecto del intervalo,
    \[
    \int_{-4}^{0}g(x)\,dx=\int_{-4}^{-2}g(x)\,dx+\int_{-2}^{0}g(x)\,dx
    \quad\Rightarrow\quad
    \int_{-2}^{0}g(x)\,dx=0-(-2)=2.
    \]

    Al invertir los límites de integración cambia el signo de la integral:
    \[
    \int_{-2}^{-4}g(x)\,dx=-\int_{-4}^{-2}g(x)\,dx=-(-2)=2.
    \]

    Sumando:
    \[
    \boxed{
    \int_{-2}^{0}g(x)\,dx+\int_{-2}^{-4}g(x)\,dx=2+2=4
    }
    \]


    \item
    \[
    \int_{-4}^{-2}\left[3g(x)-2f(x)\right]dx
    \]

    \solucion

    Por la linealidad de la integral definida:
    \[
    \int_{-4}^{-2}\left[3g(x)-2f(x)\right]dx
    =3\int_{-4}^{-2}g(x)\,dx-2\int_{-4}^{-2}f(x)\,dx
    =3(-2)-2(-5)=-6+10.
    \]

    \[
    \boxed{
    \int_{-4}^{-2}\left[3g(x)-2f(x)\right]dx=4
    }
    \]


    \item
    \[
    \int_{0}^{5}\frac52f(x)\,dx+\int_{-4}^{0}\frac12f(x)\,dx
    \]

    \solucion

    Primero obtenemos las dos integrales de $f$ que necesitamos. Por la
    aditividad respecto del intervalo:
    \[
    \int_{0}^{5}f(x)\,dx=\int_{0}^{3}f(x)\,dx+\int_{3}^{5}f(x)\,dx=7+(-4)=3,
    \]
    \[
    \int_{-4}^{0}f(x)\,dx=\int_{-4}^{5}f(x)\,dx-\int_{0}^{5}f(x)\,dx=10-3=7.
    \]

    Entonces
    \[
    \int_{0}^{5}\frac52f(x)\,dx+\int_{-4}^{0}\frac12f(x)\,dx
    =\frac52\cdot3+\frac12\cdot7=\frac{15}{2}+\frac72=\frac{22}{2}.
    \]

    \[
    \boxed{
    \int_{0}^{5}\frac52f(x)\,dx+\int_{-4}^{0}\frac12f(x)\,dx=11
    }
    \]

\end{enumerate}

\volverindice
\newpage

% ======================== EJERCICIO 10 ========================

\section{Cambio de variable en la integral definida}

Utilizando propiedades de la integral definida determine los valores de
$A$, $B$ y $C$, de modo que se satisfaga:
\[
\int_{0}^{2}x\,f(1-x^2)\,dx=C\int_{A}^{B}f(u)\,du.
\]

\solucion

Hacemos el cambio de variable sugerido por el argumento de $f$:
\[
u=1-x^2,
\qquad
du=-2x\,dx
\quad\Rightarrow\quad
x\,dx=-\frac12\,du.
\]

Al ser una integral definida, también cambiamos los límites de
integración:
\[
x=0\ \Rightarrow\ u=1-0^2=1,
\qquad
x=2\ \Rightarrow\ u=1-2^2=-3.
\]

Sustituyendo:
\[
\int_{0}^{2}x\,f(1-x^2)\,dx
=\int_{1}^{-3}f(u)\left(-\frac12\right)du
=-\frac12\int_{1}^{-3}f(u)\,du.
\]

Invirtiendo los límites de integración (lo que cambia el signo):
\[
\int_{0}^{2}x\,f(1-x^2)\,dx=\frac12\int_{-3}^{1}f(u)\,du.
\]

\[
\boxed{
A=-3,\qquad B=1,\qquad C=\frac12
}
\]

Observación: también es válida la terna equivalente $A=1$, $B=-3$,
$C=-\frac12$, que corresponde a la expresión antes de invertir los
límites.

\volverindice
\newpage

% ======================== EJERCICIO 11 ========================

\section{Área de regiones planas}

Encuentre el área de la región sombreada.

\begin{enumerate}[label=\alph*.-]

    % ----------------------------------------------------------
    % 11 a)
    % ----------------------------------------------------------

    \item Región entre $y=5x-x^2$ e $y=x$.

    \begin{center}
    \begin{tikzpicture}
        \begin{axis}[
            width=9cm, height=7cm,
            axis lines=middle,
            xlabel={$x$}, ylabel={$y$},
            xmin=-0.8, xmax=5.8,
            ymin=-1.5, ymax=7.5,
            samples=100,
            axis line style={-Latex},
        ]
            \addplot[name path=f, thick, magenta, domain=-0.3:5.3] {5*x-x^2};
            \addplot[name path=g, thick, cyan!70!blue, domain=-0.5:5.5] {x};
            \addplot[blue!15] fill between[of=f and g, soft clip={domain=0:4}];
            \addplot[only marks, mark=*, mark size=2pt, black]
                coordinates {(0,0) (4,4)};
            \node[anchor=west] at (axis cs:4.1,4) {$(4,4)$};
            \node[magenta] at (axis cs:3.6,7) {$y=5x-x^2$};
            \node[cyan!70!blue] at (axis cs:4.8,5.8) {$y=x$};
        \end{axis}
    \end{tikzpicture}
    \end{center}

    \solucion

    Puntos de intersección:
    \[
    5x-x^2=x
    \quad\Rightarrow\quad
    x^2-4x=0
    \quad\Rightarrow\quad
    x(x-4)=0
    \quad\Rightarrow\quad
    x=0,\ x=4.
    \]

    En $[0,4]$ la parábola está por encima de la recta (por ejemplo, en
    $x=1$: $4>1$). Entonces
    \[
    A=\int_{0}^{4}\left[(5x-x^2)-x\right]dx
    =\int_{0}^{4}\left(4x-x^2\right)dx
    =\left[2x^2-\frac{x^3}{3}\right]_{0}^{4}
    =32-\frac{64}{3}.
    \]

    \[
    \boxed{
    A=\frac{32}{3}\approx10{,}67
    }
    \]


    \newpage
    % ----------------------------------------------------------
    % 11 b)
    % ----------------------------------------------------------

    \item Región entre $y=\sqrt{x+2}$, $y=\dfrac{1}{x+1}$, el eje $y$
    y la recta $x=2$.

    \begin{center}
    \begin{tikzpicture}
        \begin{axis}[
            width=9cm, height=7cm,
            axis lines=middle,
            xlabel={$x$}, ylabel={$y$},
            xmin=-1.2, xmax=3.3,
            ymin=-0.5, ymax=3.3,
            samples=150,
            axis line style={-Latex},
        ]
            \addplot[name path=f, thick, magenta, domain=-2:3] {sqrt(x+2)};
            \addplot[name path=g, thick, cyan!70!blue, domain=-0.65:3] {1/(x+1)};
            \addplot[blue!15] fill between[of=f and g, soft clip={domain=0:2}];
            \draw[thick] (axis cs:2,0) -- (axis cs:2,2);
            \node[magenta] at (axis cs:1.2,2.6) {$y=\sqrt{x+2}$};
            \node[cyan!70!blue] at (axis cs:2.65,0.8) {$y=\dfrac{1}{x+1}$};
            \node[anchor=west] at (axis cs:2.05,1.3) {$x=2$};
        \end{axis}
    \end{tikzpicture}
    \end{center}

    \solucion

    En $[0,2]$ la curva $y=\sqrt{x+2}$ está por encima de
    $y=\frac{1}{x+1}$ (en $x=0$: $\sqrt2>1$, y luego la primera crece
    mientras la segunda decrece). Entonces
    \[
    A=\int_{0}^{2}\left[\sqrt{x+2}-\frac{1}{x+1}\right]dx
    =\left[\frac23(x+2)^{3/2}-\ln|x+1|\right]_{0}^{2}.
    \]

    Evaluando:
    \[
    A=\left(\frac23\cdot4^{3/2}-\ln3\right)-\left(\frac23\cdot2^{3/2}-\ln1\right)
    =\frac{16}{3}-\ln3-\frac{4\sqrt2}{3}.
    \]

    \[
    \boxed{
    A=\frac{16-4\sqrt2}{3}-\ln3\approx2{,}35
    }
    \]


    \newpage
    % ----------------------------------------------------------
    % 11 c)
    % ----------------------------------------------------------

    \item Región entre $x=e^y$, $x=y^2-2$, $y=-1$ e $y=1$.

    \begin{center}
    \begin{tikzpicture}
        \begin{axis}[
            width=9cm, height=7cm,
            axis lines=middle,
            xlabel={$x$}, ylabel={$y$},
            xmin=-2.6, xmax=4.2,
            ymin=-2.2, ymax=2.2,
            samples=150,
            axis line style={-Latex},
        ]
            \fill[blue!15]
                plot[domain=-1:1, variable=\y, samples=60] ({exp(\y)},{\y})
                -- plot[domain=1:-1, variable=\y, samples=60] ({\y*\y-2},{\y})
                -- cycle;
            \addplot[thick, magenta, variable=\t, domain=-2.2:1.4]
                ({exp(t)},{t});
            \addplot[thick, cyan!70!blue, variable=\t, domain=-2.1:2.1]
                ({t^2-2},{t});
            \draw[thick] (axis cs:-1,1) -- (axis cs:2.718,1);
            \draw[thick] (axis cs:-1,-1) -- (axis cs:0.368,-1);
            \node[magenta] at (axis cs:3.5,0.6) {$x=e^y$};
            \node[cyan!70!blue] at (axis cs:-1.4,1.8) {$x=y^2-2$};
            \node[anchor=south] at (axis cs:3,1) {$y=1$};
            \node[anchor=north] at (axis cs:1.1,-1) {$y=-1$};
        \end{axis}
    \end{tikzpicture}
    \end{center}

    \solucion

    Conviene integrar respecto de $y$, ya que las curvas están dadas
    como funciones de $y$. Para $-1\le y\le1$ la curva de la derecha es
    $x=e^y$ y la de la izquierda es $x=y^2-2$. Entonces
    \[
    A=\int_{-1}^{1}\left[e^y-(y^2-2)\right]dy
    =\int_{-1}^{1}\left(e^y-y^2+2\right)dy
    =\left[e^y-\frac{y^3}{3}+2y\right]_{-1}^{1}.
    \]

    Evaluando:
    \[
    A=\left(e-\frac13+2\right)-\left(e^{-1}+\frac13-2\right)
    =e-\frac1e+4-\frac23.
    \]

    \[
    \boxed{
    A=e-\frac1e+\frac{10}{3}\approx5{,}68
    }
    \]


    \newpage
    % ----------------------------------------------------------
    % 11 d)
    % ----------------------------------------------------------

    \item Región entre $x=y^2-4y$ y $x=2y-y^2$.

    \begin{center}
    \begin{tikzpicture}
        \begin{axis}[
            width=9cm, height=7cm,
            axis lines=middle,
            xlabel={$x$}, ylabel={$y$},
            xmin=-5.5, xmax=2.5,
            ymin=-1, ymax=4.5,
            samples=150,
            axis line style={-Latex},
        ]
            \addplot[name path=f, thick, magenta, variable=\t, domain=-0.6:3.7]
                ({2*t-t^2},{t});
            \addplot[name path=g, thick, cyan!70!blue, variable=\t, domain=-0.5:4.3]
                ({t^2-4*t},{t});
            \addplot[blue!15] fill between[of=f and g,
                soft clip={(axis cs:-6,0) rectangle (axis cs:3,3)}];
            \addplot[only marks, mark=*, mark size=2pt, black]
                coordinates {(0,0) (-3,3)};
            \node[anchor=east] at (axis cs:-3.1,3.1) {$(-3,3)$};
            \node[cyan!70!blue] at (axis cs:-2.4,4.1) {$x=y^2-4y$};
            \node[magenta] at (axis cs:-3.6,0.5) {$x=2y-y^2$};
        \end{axis}
    \end{tikzpicture}
    \end{center}

    \solucion

    Integramos respecto de $y$. Puntos de intersección:
    \[
    y^2-4y=2y-y^2
    \quad\Rightarrow\quad
    2y^2-6y=0
    \quad\Rightarrow\quad
    2y(y-3)=0
    \quad\Rightarrow\quad
    y=0,\ y=3,
    \]
    que corresponden a los puntos $(0,0)$ y $(-3,3)$.

    En $[0,3]$ la curva de la derecha es $x=2y-y^2$ (en $y=1$: $1>-3$).
    Entonces
    \[
    A=\int_{0}^{3}\left[(2y-y^2)-(y^2-4y)\right]dy
    =\int_{0}^{3}\left(6y-2y^2\right)dy
    =\left[3y^2-\frac23y^3\right]_{0}^{3}
    =27-18.
    \]

    \[
    \boxed{
    A=9
    }
    \]


    \newpage
    % ----------------------------------------------------------
    % 11 e)
    % ----------------------------------------------------------

    \item Región entre $y=\frac12\sec^2t$ e $y=-4\sen^2t$, para
    $-\frac\pi3\le t\le\frac\pi3$.

    \begin{center}
    \begin{tikzpicture}
        \begin{axis}[
            width=9cm, height=7.5cm,
            axis lines=middle,
            xlabel={$t$}, ylabel={$y$},
            xtick={-1.0471975512, 1.0471975512},
            xticklabels={$-\frac{\pi}{3}$, $\frac{\pi}{3}$},
            ytick={-4,-2,1,2},
            xmin=-1.6, xmax=1.6,
            ymin=-4.6, ymax=3.2,
            samples=150,
            restrict y to domain=-5:4,
            axis line style={-Latex},
        ]
            \addplot[name path=f, thick, magenta, domain=-1.2:1.2]
                {0.5/(cos(deg(x)))^2};
            \addplot[name path=g, thick, cyan!70!blue, domain=-1.5:1.5]
                {-4*(sin(deg(x)))^2};
            \addplot[blue!15] fill between[of=f and g,
                soft clip={domain=-1.0471975512:1.0471975512}];
            \draw[dashed] (axis cs:-1.0471975512,-3) -- (axis cs:-1.0471975512,2);
            \draw[dashed] (axis cs:1.0471975512,-3) -- (axis cs:1.0471975512,2);
            \node[magenta] at (axis cs:0.55,2.1) {$y=\frac12\sec^2t$};
            \node[cyan!70!blue] at (axis cs:0.6,-4.2) {$y=-4\sen^2t$};
        \end{axis}
    \end{tikzpicture}
    \end{center}

    \solucion

    En $\left[-\frac\pi3,\frac\pi3\right]$ se tiene
    $\frac12\sec^2t>0\ge-4\sen^2t$, de modo que
    \[
    A=\int_{-\pi/3}^{\pi/3}\left[\frac12\sec^2t+4\sen^2t\right]dt.
    \]

    Como el integrando es una función par, podemos integrar en
    $\left[0,\frac\pi3\right]$ y duplicar:
    \[
    A=2\int_{0}^{\pi/3}\left[\frac12\sec^2t+4\sen^2t\right]dt.
    \]

    Calculamos cada parte. La primera es inmediata:
    \[
    \int_{0}^{\pi/3}\frac12\sec^2t\,dt=\frac12\Big[\tan t\Big]_{0}^{\pi/3}
    =\frac{\sqrt3}{2}.
    \]

    Para la segunda usamos $\sen^2t=\dfrac{1-\cos(2t)}{2}$:
    \[
    \int_{0}^{\pi/3}4\sen^2t\,dt
    =\int_{0}^{\pi/3}\left(2-2\cos(2t)\right)dt
    =\Big[2t-\sen(2t)\Big]_{0}^{\pi/3}
    =\frac{2\pi}{3}-\frac{\sqrt3}{2}.
    \]

    Sumando: $\dfrac{\sqrt3}{2}+\dfrac{2\pi}{3}-\dfrac{\sqrt3}{2}=\dfrac{2\pi}{3}$.
    Por lo tanto
    \[
    \boxed{
    A=2\cdot\frac{2\pi}{3}=\frac{4\pi}{3}\approx4{,}19
    }
    \]


    \newpage
    % ----------------------------------------------------------
    % 11 f)
    % ----------------------------------------------------------

    \item (*) Región entre $y=-x^2+3x$ e $y=2x^3-x^2-5x$.

    \begin{center}
    \begin{tikzpicture}
        \begin{axis}[
            width=9cm, height=7.5cm,
            axis lines=middle,
            xlabel={$x$}, ylabel={$y$},
            xmin=-2.6, xmax=2.8,
            ymin=-11.5, ymax=4.5,
            ytick={-5},
            samples=150,
            restrict y to domain=-13:6,
            axis line style={-Latex},
        ]
            \addplot[name path=f, thick, cyan!70!blue, domain=-2.1:3]
                {-x^2+3*x};
            \addplot[name path=g, thick, black, domain=-2.05:2.4]
                {2*x^3-x^2-5*x};
            \addplot[blue!15] fill between[of=f and g, soft clip={domain=-2:2}];
            \addplot[only marks, mark=*, mark size=2pt, black]
                coordinates {(-2,-10) (0,0) (2,2)};
            \node[anchor=west] at (axis cs:-1.95,-10.3) {$(-2,-10)$};
            \node[anchor=west] at (axis cs:2.05,1.6) {$(2,2)$};
            \node[cyan!70!blue] at (axis cs:1.45,3.9) {$y=-x^2+3x$};
            \node[black] at (axis cs:1.5,-6) {$y=2x^3-x^2-5x$};
        \end{axis}
    \end{tikzpicture}
    \end{center}

    \solucion

    Puntos de intersección:
    \[
    2x^3-x^2-5x=-x^2+3x
    \quad\Rightarrow\quad
    2x^3-8x=0,
    \]
    \[
    2x(x-2)(x+2)=0
    \quad\Rightarrow\quad
    x=-2,\ 0,\ 2.
    \]

    Las curvas se cruzan en $x=0$, por lo que el área se divide en dos
    partes:
    \begin{itemize}
        \item En $[-2,0]$ la cúbica está arriba (en $x=-1$: $2>-4$).
        \item En $[0,2]$ la parábola está arriba (en $x=1$: $2>-4$).
    \end{itemize}

    La diferencia entre ambas es $(2x^3-x^2-5x)-(-x^2+3x)=2x^3-8x$. Entonces
    \[
    A=\int_{-2}^{0}\left(2x^3-8x\right)dx+\int_{0}^{2}\left(8x-2x^3\right)dx.
    \]

    \[
    \int_{-2}^{0}\left(2x^3-8x\right)dx
    =\left[\frac{x^4}{2}-4x^2\right]_{-2}^{0}=0-(8-16)=8,
    \]
    \[
    \int_{0}^{2}\left(8x-2x^3\right)dx
    =\left[4x^2-\frac{x^4}{2}\right]_{0}^{2}=16-8=8.
    \]

    \[
    \boxed{
    A=8+8=16
    }
    \]

\end{enumerate}

\volverindice
\newpage

% ======================== EJERCICIO 12 ========================

\section{Integrales impropias}

Determine si cada una de las siguientes integrales es convergente o
divergente. Evalúe las que sean convergentes.

\begin{enumerate}[label=\alph*.-]

    \item
    \[
    \int_{1}^{\infty}\frac{\ln y}{y^3}\,dy
    \]

    \solucion

    Es una integral impropia de primera especie (intervalo no acotado):
    \[
    \int_{1}^{\infty}\frac{\ln y}{y^3}\,dy
    =\lim_{b\to\infty}\int_{1}^{b}\frac{\ln y}{y^3}\,dy.
    \]

    Calculamos la primitiva por partes, eligiendo
    \[
    u=\ln y,\qquad dv=y^{-3}\,dy
    \quad\Rightarrow\quad
    du=\frac1y\,dy,\qquad v=-\frac{1}{2y^2}.
    \]
    \[
    \int\frac{\ln y}{y^3}\,dy
    =-\frac{\ln y}{2y^2}+\int\frac{1}{2y^3}\,dy
    =-\frac{\ln y}{2y^2}-\frac{1}{4y^2}+C.
    \]

    Entonces
    \[
    \int_{1}^{b}\frac{\ln y}{y^3}\,dy
    =\left(-\frac{\ln b}{2b^2}-\frac{1}{4b^2}\right)-\left(0-\frac14\right).
    \]

    Por la regla de L'Hôpital,
    $\displaystyle\lim_{b\to\infty}\frac{\ln b}{b^2}=\lim_{b\to\infty}\frac{1/b}{2b}=0$,
    y también $\dfrac{1}{4b^2}\to0$. Por lo tanto el límite existe y vale
    $\frac14$.

    \[
    \boxed{
    \int_{1}^{\infty}\frac{\ln y}{y^3}\,dy=\frac14
    \quad\text{(convergente)}
    }
    \]


    \item
    \[
    \int_{-\infty}^{0}t\,e^{2t}\,dt
    \]

    \solucion

    \[
    \int_{-\infty}^{0}t\,e^{2t}\,dt=\lim_{a\to-\infty}\int_{a}^{0}t\,e^{2t}\,dt.
    \]

    Por partes, con $u=t$, $dv=e^{2t}\,dt$ ($du=dt$, $v=\frac12e^{2t}$):
    \[
    \int t\,e^{2t}\,dt=\frac{t}{2}e^{2t}-\frac12\int e^{2t}\,dt
    =\frac{t}{2}e^{2t}-\frac14e^{2t}+C.
    \]

    Entonces
    \[
    \int_{a}^{0}t\,e^{2t}\,dt
    =\left(0-\frac14\right)-\left(\frac{a}{2}e^{2a}-\frac14e^{2a}\right).
    \]

    Cuando $a\to-\infty$: $e^{2a}\to0$ y, por L'Hôpital,
    \[
    \lim_{a\to-\infty}a\,e^{2a}=\lim_{a\to-\infty}\frac{a}{e^{-2a}}
    =\lim_{a\to-\infty}\frac{1}{-2e^{-2a}}=0.
    \]

    \[
    \boxed{
    \int_{-\infty}^{0}t\,e^{2t}\,dt=-\frac14
    \quad\text{(convergente)}
    }
    \]


    \item
    \[
    \int_{-\infty}^{\infty}z\,e^{-z^2}\,dz
    \]

    \solucion

    Separamos en $z=0$; la integral converge sólo si convergen ambas
    partes:
    \[
    \int_{-\infty}^{\infty}z\,e^{-z^2}\,dz
    =\int_{-\infty}^{0}z\,e^{-z^2}\,dz+\int_{0}^{\infty}z\,e^{-z^2}\,dz.
    \]

    Con $w=-z^2$, $dw=-2z\,dz$, la primitiva es
    $\displaystyle\int z\,e^{-z^2}\,dz=-\frac12e^{-z^2}+C$. Entonces
    \[
    \int_{0}^{\infty}z\,e^{-z^2}\,dz
    =\lim_{b\to\infty}\left[-\frac12e^{-z^2}\right]_{0}^{b}
    =\lim_{b\to\infty}\left(-\frac12e^{-b^2}+\frac12\right)=\frac12,
    \]
    \[
    \int_{-\infty}^{0}z\,e^{-z^2}\,dz
    =\lim_{a\to-\infty}\left[-\frac12e^{-z^2}\right]_{a}^{0}
    =\lim_{a\to-\infty}\left(-\frac12+\frac12e^{-a^2}\right)=-\frac12.
    \]

    Ambas convergen, así que
    \[
    \boxed{
    \int_{-\infty}^{\infty}z\,e^{-z^2}\,dz=-\frac12+\frac12=0
    \quad\text{(convergente)}
    }
    \]


    \item
    \[
    \int_{-\infty}^{\infty}\frac{e^x}{1+e^{2x}}\,dx
    \]

    \solucion

    Con la sustitución $u=e^x$, $du=e^x\,dx$:
    \[
    \int\frac{e^x}{1+e^{2x}}\,dx=\int\frac{du}{1+u^2}=\arctan(e^x)+C.
    \]

    Separamos en $x=0$:
    \[
    \int_{0}^{\infty}\frac{e^x}{1+e^{2x}}\,dx
    =\lim_{b\to\infty}\Big[\arctan(e^x)\Big]_{0}^{b}
    =\frac\pi2-\frac\pi4=\frac\pi4,
    \]
    ya que $e^b\to\infty$ y $\arctan(e^b)\to\frac\pi2$.
    \[
    \int_{-\infty}^{0}\frac{e^x}{1+e^{2x}}\,dx
    =\lim_{a\to-\infty}\Big[\arctan(e^x)\Big]_{a}^{0}
    =\frac\pi4-0=\frac\pi4,
    \]
    ya que $e^a\to0$ y $\arctan(e^a)\to0$.

    \[
    \boxed{
    \int_{-\infty}^{\infty}\frac{e^x}{1+e^{2x}}\,dx=\frac\pi4+\frac\pi4=\frac\pi2
    \quad\text{(convergente)}
    }
    \]


    \item
    \[
    \int_{0}^{1}\frac{x+1}{\sqrt{x^2+2x}}\,dx
    \]

    \solucion

    Es una integral impropia de segunda especie: en $x=0$ el denominador
    se anula ($\sqrt{0^2+2\cdot0}=0$) y el integrando no está acotado.
    \[
    \int_{0}^{1}\frac{x+1}{\sqrt{x^2+2x}}\,dx
    =\lim_{a\to0^+}\int_{a}^{1}\frac{x+1}{\sqrt{x^2+2x}}\,dx.
    \]

    Con $u=x^2+2x$, $du=2(x+1)\,dx$:
    \[
    \int\frac{x+1}{\sqrt{x^2+2x}}\,dx=\frac12\int u^{-1/2}\,du=\sqrt u+C
    =\sqrt{x^2+2x}+C.
    \]

    Entonces
    \[
    \lim_{a\to0^+}\left[\sqrt{x^2+2x}\right]_{a}^{1}
    =\lim_{a\to0^+}\left(\sqrt3-\sqrt{a^2+2a}\right)=\sqrt3.
    \]

    \[
    \boxed{
    \int_{0}^{1}\frac{x+1}{\sqrt{x^2+2x}}\,dx=\sqrt3
    \quad\text{(convergente)}
    }
    \]


    \item
    \[
    \int_{0}^{4}\frac{dx}{\sqrt{16-x^2}}
    \]

    \solucion

    El integrando no está acotado en $x=4$, por lo que
    \[
    \int_{0}^{4}\frac{dx}{\sqrt{16-x^2}}
    =\lim_{b\to4^-}\int_{0}^{b}\frac{dx}{\sqrt{16-x^2}}.
    \]

    La primitiva es inmediata:
    $\displaystyle\int\frac{dx}{\sqrt{16-x^2}}=\operatorname{arcsen}\left(\frac x4\right)+C$.
    Entonces
    \[
    \lim_{b\to4^-}\left[\operatorname{arcsen}\left(\frac x4\right)\right]_{0}^{b}
    =\operatorname{arcsen}(1)-\operatorname{arcsen}(0)=\frac\pi2.
    \]

    \[
    \boxed{
    \int_{0}^{4}\frac{dx}{\sqrt{16-x^2}}=\frac\pi2
    \quad\text{(convergente)}
    }
    \]


    \item (*)
    \[
    \int_{-1}^{1}\frac{1}{x^2}\,dx
    \]

    \solucion

    El integrando no está definido en $x=0$, que es un punto interior del
    intervalo. Separamos:
    \[
    \int_{-1}^{1}\frac{dx}{x^2}=\int_{-1}^{0}\frac{dx}{x^2}+\int_{0}^{1}\frac{dx}{x^2}.
    \]

    Analizamos la segunda:
    \[
    \int_{0}^{1}\frac{dx}{x^2}
    =\lim_{a\to0^+}\left[-\frac1x\right]_{a}^{1}
    =\lim_{a\to0^+}\left(-1+\frac1a\right)=+\infty.
    \]

    Como una de las partes diverge, la integral completa diverge.

    \[
    \boxed{
    \int_{-1}^{1}\frac{1}{x^2}\,dx\ \text{ es divergente}
    }
    \]

    Observación: si se ignorara la discontinuidad y se aplicara
    directamente la regla de Barrow, se obtendría
    $\left[-\frac1x\right]_{-1}^{1}=-2$, un resultado absurdo, ya que el
    integrando es siempre positivo.


    \item (*)
    \[
    \int_{0}^{2}\frac{dx}{x^2-6x+5}
    \]

    \solucion

    Factorizamos el denominador: $x^2-6x+5=(x-1)(x-5)$. En el intervalo
    $[0,2]$ el integrando no está definido en $x=1$, así que
    \[
    \int_{0}^{2}\frac{dx}{x^2-6x+5}
    =\int_{0}^{1}\frac{dx}{(x-1)(x-5)}+\int_{1}^{2}\frac{dx}{(x-1)(x-5)}.
    \]

    Por fracciones parciales,
    $\dfrac{1}{(x-1)(x-5)}=\dfrac{A}{x-1}+\dfrac{B}{x-5}$ con
    $A=-\frac14$ y $B=\frac14$, de modo que
    \[
    \int\frac{dx}{(x-1)(x-5)}=\frac14\ln\left|\frac{x-5}{x-1}\right|+C.
    \]

    Analizamos la primera parte:
    \[
    \int_{0}^{1}\frac{dx}{(x-1)(x-5)}
    =\lim_{b\to1^-}\left[\frac14\ln\left|\frac{x-5}{x-1}\right|\right]_{0}^{b}
    =\lim_{b\to1^-}\frac14\ln\left|\frac{b-5}{b-1}\right|-\frac14\ln5=+\infty,
    \]
    ya que $\left|\frac{b-5}{b-1}\right|\to+\infty$ cuando $b\to1^-$.

    \[
    \boxed{
    \int_{0}^{2}\frac{dx}{x^2-6x+5}\ \text{ es divergente}
    }
    \]

\end{enumerate}

\volverindice
\newpage

% ======================== EJERCICIO 13 ========================

\section{Longitud de arco}

Determine la longitud de arco de las siguientes curvas.

Recordamos que la longitud de arco de una curva $y=f(x)$, con
$a\le x\le b$, es
\[
L=\int_{a}^{b}\sqrt{1+\left[f'(x)\right]^2}\,dx,
\]
y análogamente, para una curva $x=g(y)$ con $c\le y\le d$,
\[
L=\int_{c}^{d}\sqrt{1+\left[g'(y)\right]^2}\,dy.
\]

\begin{enumerate}[label=\alph*.-]

    \item $y=1+6x^{3/2}$, \quad $0\le x\le1$.

    \solucion

    Derivando: $y'=6\cdot\frac32x^{1/2}=9\sqrt x$, de donde
    $(y')^2=81x$. Entonces
    \[
    L=\int_{0}^{1}\sqrt{1+81x}\,dx.
    \]

    Con $u=1+81x$, $du=81\,dx$ (si $x=0$, $u=1$; si $x=1$, $u=82$):
    \[
    L=\frac{1}{81}\int_{1}^{82}u^{1/2}\,du
    =\frac{1}{81}\cdot\frac23\Big[u^{3/2}\Big]_{1}^{82}
    =\frac{2}{243}\left(82^{3/2}-1\right).
    \]

    \[
    \boxed{
    L=\frac{2}{243}\left(82\sqrt{82}-1\right)\approx6{,}10
    }
    \]


    \item $x=\dfrac{y^4}{8}+\dfrac{1}{4y^2}$, \quad $1\le y\le2$.

    \solucion

    Derivando respecto de $y$:
    \[
    \frac{dx}{dy}=\frac{y^3}{2}-\frac{1}{2y^3}.
    \]

    Elevando al cuadrado:
    \[
    \left(\frac{dx}{dy}\right)^2=\frac{y^6}{4}-\frac12+\frac{1}{4y^6},
    \]
    y sumando 1 se obtiene un cuadrado perfecto:
    \[
    1+\left(\frac{dx}{dy}\right)^2=\frac{y^6}{4}+\frac12+\frac{1}{4y^6}
    =\left(\frac{y^3}{2}+\frac{1}{2y^3}\right)^2.
    \]

    Como $\frac{y^3}{2}+\frac{1}{2y^3}>0$ en $[1,2]$:
    \[
    L=\int_{1}^{2}\left(\frac{y^3}{2}+\frac{1}{2y^3}\right)dy
    =\left[\frac{y^4}{8}-\frac{1}{4y^2}\right]_{1}^{2}
    =\left(2-\frac{1}{16}\right)-\left(\frac18-\frac14\right)
    =\frac{31}{16}+\frac18.
    \]

    \[
    \boxed{
    L=\frac{33}{16}
    }
    \]


    \item $x=\dfrac13\sqrt y\,(y-3)$, \quad $1\le y\le9$.

    \solucion

    Escribimos $x=\frac13\left(y^{3/2}-3y^{1/2}\right)$ y derivamos:
    \[
    \frac{dx}{dy}=\frac13\left(\frac32y^{1/2}-\frac32y^{-1/2}\right)
    =\frac12\left(y^{1/2}-y^{-1/2}\right).
    \]

    Entonces
    \[
    1+\left(\frac{dx}{dy}\right)^2
    =1+\frac14\left(y-2+\frac1y\right)
    =\frac14\left(y+2+\frac1y\right)
    =\left[\frac12\left(y^{1/2}+y^{-1/2}\right)\right]^2.
    \]

    Por lo tanto
    \[
    L=\int_{1}^{9}\frac12\left(y^{1/2}+y^{-1/2}\right)dy
    =\frac12\left[\frac23y^{3/2}+2y^{1/2}\right]_{1}^{9},
    \]
    \[
    L=\frac12\left[(18+6)-\left(\frac23+2\right)\right]
    =\frac12\cdot\frac{64}{3}.
    \]

    \[
    \boxed{
    L=\frac{32}{3}
    }
    \]


    \item $y=\ln(\cos x)$, \quad $0\le x\le\dfrac\pi3$.

    \solucion

    Derivando: $y'=\dfrac{-\sen x}{\cos x}=-\tan x$. Por la identidad
    pitagórica:
    \[
    1+(y')^2=1+\tan^2x=\sec^2x.
    \]

    Como $\sec x>0$ en $\left[0,\frac\pi3\right]$:
    \[
    L=\int_{0}^{\pi/3}\sec x\,dx
    =\Big[\ln|\sec x+\tan x|\Big]_{0}^{\pi/3}
    =\ln\left(2+\sqrt3\right)-\ln(1+0).
    \]

    \[
    \boxed{
    L=\ln\left(2+\sqrt3\right)\approx1{,}32
    }
    \]


    \item (*) $y=\dfrac14x^2-\dfrac12\ln x$, \quad $1\le x\le2$.

    \solucion

    Derivando: $y'=\dfrac x2-\dfrac{1}{2x}$. Entonces
    \[
    1+(y')^2=1+\frac{x^2}{4}-\frac12+\frac{1}{4x^2}
    =\frac{x^2}{4}+\frac12+\frac{1}{4x^2}
    =\left(\frac x2+\frac{1}{2x}\right)^2.
    \]

    Por lo tanto
    \[
    L=\int_{1}^{2}\left(\frac x2+\frac{1}{2x}\right)dx
    =\left[\frac{x^2}{4}+\frac12\ln x\right]_{1}^{2}
    =\left(1+\frac12\ln2\right)-\frac14.
    \]

    \[
    \boxed{
    L=\frac34+\frac12\ln2\approx1{,}10
    }
    \]


    \item (*) $y=1-e^{-x}$, \quad $0\le x\le2$.

    \solucion

    Derivando: $y'=e^{-x}$, así que
    \[
    L=\int_{0}^{2}\sqrt{1+e^{-2x}}\,dx.
    \]

    Sustituimos $w=e^{-x}$, $dw=-e^{-x}\,dx=-w\,dx$, es decir,
    $dx=-\dfrac{dw}{w}$. Los límites cambian a $x=0\Rightarrow w=1$ y
    $x=2\Rightarrow w=e^{-2}$:
    \[
    L=\int_{1}^{e^{-2}}\sqrt{1+w^2}\left(-\frac{dw}{w}\right)
    =\int_{e^{-2}}^{1}\frac{\sqrt{1+w^2}}{w}\,dw.
    \]

    Esta integral fue resuelta por sustitución trigonométrica en el
    ejercicio 7.c de la parte 1:
    \[
    \int\frac{\sqrt{1+w^2}}{w}\,dw
    =\sqrt{1+w^2}-\ln\left|\frac{1+\sqrt{1+w^2}}{w}\right|+C.
    \]

    Evaluando en $w=1$:
    \[
    \sqrt2-\ln\left(1+\sqrt2\right).
    \]

    Evaluando en $w=e^{-2}$, y usando que
    $\ln\left(\dfrac{1+\sqrt{1+e^{-4}}}{e^{-2}}\right)=\ln\left(1+\sqrt{1+e^{-4}}\right)+2$:
    \[
    \sqrt{1+e^{-4}}-\ln\left(1+\sqrt{1+e^{-4}}\right)-2.
    \]

    Restando:
    \[
    \boxed{
    L=2+\sqrt2-\sqrt{1+e^{-4}}+\ln\left(\frac{1+\sqrt{1+e^{-4}}}{1+\sqrt2}\right)
    \approx2{,}22
    }
    \]

\end{enumerate}

\volverindice
\newpage

% ======================== EJERCICIO 14 ========================

\section{Trabajo}

Resolver lo pedido en cada caso. Recordamos que el trabajo realizado por
una fuerza variable $F(x)$ al mover un objeto desde $x=a$ hasta $x=b$ es
\[
W=\int_{a}^{b}F(x)\,dx,
\]
y que, por la ley de Hooke, la fuerza necesaria para mantener un
resorte estirado $x$ unidades más allá de su longitud natural es
$F(x)=kx$.

\begin{enumerate}[label=\alph*.-]

    \item Una fuerza variable de $5x^{-2}\ \text{N}$ mueve un objeto a
    lo largo de una línea recta cuando está a $x$ metros del origen.
    Calcule el trabajo realizado para mover el objeto desde $x=1\ \text{m}$
    a $x=10\ \text{m}$.

    \solucion

    \[
    W=\int_{1}^{10}5x^{-2}\,dx
    =5\left[-\frac1x\right]_{1}^{10}
    =5\left(-\frac{1}{10}+1\right)
    =5\cdot\frac{9}{10}.
    \]

    \[
    \boxed{
    W=4{,}5\ \text{J}
    }
    \]


    \item Cuando una partícula se localiza a una distancia de $x$ metros
    desde el origen, una fuerza de $\cos(\pi x/3)$ newtons actúa sobre
    ella. ¿Cuánto trabajo se realiza al mover la partícula desde $x=1$
    hasta $x=2$? Interprete su respuesta considerando el trabajo
    realizado desde $x=1$ hasta $x=1{,}5$ y desde $x=1{,}5$ hasta $x=2$.

    \solucion

    La primitiva de $\cos\left(\frac{\pi x}{3}\right)$ es
    $\frac3\pi\sen\left(\frac{\pi x}{3}\right)$. Entonces
    \[
    W=\int_{1}^{2}\cos\left(\frac{\pi x}{3}\right)dx
    =\frac3\pi\left[\sen\left(\frac{\pi x}{3}\right)\right]_{1}^{2},
    \]
    \[
    W=\frac3\pi\left(\sen\frac{2\pi}{3}-\sen\frac\pi3\right)
    =\frac3\pi\left(\frac{\sqrt3}{2}-\frac{\sqrt3}{2}\right).
    \]

    \[
    \boxed{
    W=0\ \text{J}
    }
    \]

    \textbf{Interpretación.} Calculamos el trabajo en cada mitad del
    recorrido:
    \[
    W_1=\int_{1}^{1{,}5}\cos\left(\frac{\pi x}{3}\right)dx
    =\frac3\pi\left(\sen\frac\pi2-\sen\frac\pi3\right)
    =\frac3\pi\left(1-\frac{\sqrt3}{2}\right)\approx0{,}128\ \text{J},
    \]
    \[
    W_2=\int_{1{,}5}^{2}\cos\left(\frac{\pi x}{3}\right)dx
    =\frac3\pi\left(\sen\frac{2\pi}{3}-\sen\frac\pi2\right)
    =\frac3\pi\left(\frac{\sqrt3}{2}-1\right)\approx-0{,}128\ \text{J}.
    \]

    La fuerza $\cos\left(\frac{\pi x}{3}\right)$ se anula en $x=1{,}5$:
    entre $x=1$ y $x=1{,}5$ es positiva (actúa en el sentido del
    movimiento y realiza trabajo positivo), mientras que entre $x=1{,}5$
    y $x=2$ es negativa (se opone al movimiento). Por la simetría del
    coseno respecto de $x=1{,}5$, ambos trabajos tienen el mismo valor
    absoluto y signos opuestos, por lo que se cancelan y el trabajo
    neto es nulo.


    \item Un resorte tiene una longitud natural de $20\ \text{cm}$. Si
    se requiere una fuerza de $25\ \text{N}$ para mantenerlo estirado a
    una longitud de $30\ \text{cm}$, ¿cuánto trabajo se requiere para
    estirarlo desde $20$ hasta $25\ \text{cm}$?

    \solucion

    Al estar a $30\ \text{cm}$, el resorte está estirado
    $30-20=10\ \text{cm}=0{,}1\ \text{m}$. Por la ley de Hooke:
    \[
    25=k\cdot0{,}1
    \quad\Rightarrow\quad
    k=250\ \text{N/m}.
    \]

    Estirarlo desde $20$ hasta $25\ \text{cm}$ equivale a llevar el
    estiramiento desde $x=0$ hasta $x=0{,}05\ \text{m}$:
    \[
    W=\int_{0}^{0{,}05}250x\,dx
    =125\Big[x^2\Big]_{0}^{0{,}05}
    =125\cdot0{,}0025.
    \]

    \[
    \boxed{
    W=0{,}3125\ \text{J}
    }
    \]


    \item Suponga que se necesitan $2\ \text{J}$ de trabajo para estirar
    un resorte desde su longitud natural de $30\ \text{cm}$ hasta una
    longitud de $42\ \text{cm}$.

    \solucion

    Primero hallamos la constante $k$. Estirar de $30$ a $42\ \text{cm}$
    es llevar el estiramiento de $0$ a $0{,}12\ \text{m}$:
    \[
    2=\int_{0}^{0{,}12}kx\,dx=\frac k2(0{,}12)^2=0{,}0072\,k,
    \]
    \[
    k=\frac{2}{0{,}0072}=\frac{2500}{9}\ \text{N/m}\approx277{,}8\ \text{N/m}.
    \]

    \begin{enumerate}[label=\roman*)]

        \item ¿Cuánto trabajo se requiere para estirarlo desde $35$
        hasta $40\ \text{cm}$?

        Corresponde a estiramientos desde $x=0{,}05$ hasta
        $x=0{,}10\ \text{m}$:
        \[
        W=\int_{0{,}05}^{0{,}10}\frac{2500}{9}x\,dx
        =\frac{1250}{9}\Big[x^2\Big]_{0{,}05}^{0{,}10}
        =\frac{1250}{9}\left(0{,}01-0{,}0025\right)
        =\frac{1250}{9}\cdot0{,}0075.
        \]

        \[
        \boxed{
        W=\frac{25}{24}\ \text{J}\approx1{,}04\ \text{J}
        }
        \]

        \item (*) ¿Cuánto más allá de su longitud natural, una fuerza
        de $30\ \text{N}$ mantendrá el resorte estirado?

        Por la ley de Hooke:
        \[
        30=\frac{2500}{9}\,x
        \quad\Rightarrow\quad
        x=\frac{270}{2500}=0{,}108\ \text{m}.
        \]

        \[
        \boxed{
        x=10{,}8\ \text{cm}
        }
        \]

    \end{enumerate}

\end{enumerate}

\volverindice
\newpage

% ======================== EJERCICIO 15 ========================

\section{Masa y centro de masa de una barra}

Calcule la masa total y el centro de masa de la barra indicada.

Si una barra ocupa el intervalo $[a,b]$ y su densidad lineal es
$\rho(x)$, su masa total y su centro de masa son
\[
M=\int_{a}^{b}\rho(x)\,dx,
\qquad
\bar x=\frac{1}{M}\int_{a}^{b}x\,\rho(x)\,dx.
\]

En todos los casos ubicamos el extremo de referencia de la barra en
$x=0$.

\begin{enumerate}[label=\alph*.-]

    \item La longitud de una barra es de $6\ \text{m}$ y la densidad
    lineal de la barra en un punto que está a $x$ metros de un extremo es
    $2x+3\ \text{kg/m}$.

    \solucion

    \[
    M=\int_{0}^{6}(2x+3)\,dx=\Big[x^2+3x\Big]_{0}^{6}=36+18=54\ \text{kg}.
    \]

    Momento respecto del origen:
    \[
    \int_{0}^{6}x(2x+3)\,dx=\int_{0}^{6}\left(2x^2+3x\right)dx
    =\left[\frac23x^3+\frac32x^2\right]_{0}^{6}=144+54=198.
    \]

    \[
    \bar x=\frac{198}{54}=\frac{11}{3}.
    \]

    \[
    \boxed{
    M=54\ \text{kg},
    \qquad
    \bar x=\frac{11}{3}\ \text{m}\approx3{,}67\ \text{m del extremo}
    }
    \]

    Como era de esperar, el centro de masa está desplazado hacia el
    extremo más denso (más allá del punto medio, $x=3$).


    \item La longitud de la barra es de $9\ \text{pulg}$ y la densidad
    lineal de la barra en un punto que está a $x$ pulgadas de un extremo
    es $4x+1\ \text{slugs/pulg}$.

    \solucion

    \[
    M=\int_{0}^{9}(4x+1)\,dx=\Big[2x^2+x\Big]_{0}^{9}=162+9=171\ \text{slugs}.
    \]

    Momento respecto del origen:
    \[
    \int_{0}^{9}x(4x+1)\,dx=\int_{0}^{9}\left(4x^2+x\right)dx
    =\left[\frac43x^3+\frac12x^2\right]_{0}^{9}=972+\frac{81}{2}=\frac{2025}{2}.
    \]

    \[
    \bar x=\frac{2025/2}{171}=\frac{2025}{342}=\frac{225}{38}.
    \]

    \[
    \boxed{
    M=171\ \text{slugs},
    \qquad
    \bar x=\frac{225}{38}\ \text{pulg}\approx5{,}92\ \text{pulg del extremo}
    }
    \]


    \item La longitud de la barra es de $12\ \text{cm}$ y la medida de
    la densidad lineal de la barra en un punto es una función lineal de
    la medida de la distancia del punto al extremo izquierdo de la barra.
    La densidad lineal en el extremo izquierdo es $3\ \text{g/cm}$ y en
    el extremo derecho es $4\ \text{g/cm}$.

    \solucion

    Ubicamos el extremo izquierdo en $x=0$. Como la densidad es lineal,
    $\rho(x)=mx+n$, con $\rho(0)=3$ y $\rho(12)=4$:
    \[
    n=3,
    \qquad
    m=\frac{4-3}{12-0}=\frac{1}{12}
    \quad\Rightarrow\quad
    \rho(x)=3+\frac{x}{12}.
    \]

    Masa total:
    \[
    M=\int_{0}^{12}\left(3+\frac{x}{12}\right)dx
    =\left[3x+\frac{x^2}{24}\right]_{0}^{12}=36+6=42\ \text{g}.
    \]

    Momento respecto del extremo izquierdo:
    \[
    \int_{0}^{12}x\left(3+\frac{x}{12}\right)dx
    =\left[\frac32x^2+\frac{x^3}{36}\right]_{0}^{12}=216+48=264.
    \]

    \[
    \bar x=\frac{264}{42}=\frac{44}{7}.
    \]

    \[
    \boxed{
    M=42\ \text{g},
    \qquad
    \bar x=\frac{44}{7}\ \text{cm}\approx6{,}29\ \text{cm del extremo izquierdo}
    }
    \]


    \item (*) La medida de la densidad lineal en cualquier punto de una
    barra de $6\ \text{m}$ de longitud varía directamente como la
    distancia del punto a un punto externo situado a $4\ \text{m}$ de un
    extremo, donde la densidad es $3\ \text{kg/m}$.

    \solucion

    Colocamos el punto externo en el origen, sobre la recta de la barra.
    Entonces el extremo más cercano está en $x=4$ y el más lejano en
    $x=10$. Como la densidad varía directamente con la distancia al
    origen, $\rho(x)=kx$. En el extremo cercano ($x=4$) la densidad es
    $3\ \text{kg/m}$:
    \[
    3=4k
    \quad\Rightarrow\quad
    k=\frac34
    \quad\Rightarrow\quad
    \rho(x)=\frac34x.
    \]

    Masa total:
    \[
    M=\int_{4}^{10}\frac34x\,dx=\frac38\Big[x^2\Big]_{4}^{10}
    =\frac38(100-16)=\frac{63}{2}=31{,}5\ \text{kg}.
    \]

    Momento respecto del punto externo:
    \[
    \int_{4}^{10}\frac34x^2\,dx=\frac14\Big[x^3\Big]_{4}^{10}
    =\frac14(1000-64)=234.
    \]

    \[
    \bar x=\frac{234}{63/2}=\frac{468}{63}=\frac{52}{7}\ \text{m}
    \]
    medidos desde el punto externo. Respecto del extremo más cercano a
    dicho punto: $\frac{52}{7}-4=\frac{24}{7}$.

    \[
    \boxed{
    M=31{,}5\ \text{kg},
    \qquad
    \bar x=\frac{24}{7}\ \text{m}\approx3{,}43\ \text{m del extremo más cercano al punto externo}
    }
    \]

\end{enumerate}

\volverindice
\newpage

% ======================== EJERCICIO 16 ========================

\section{Verificación de soluciones de ecuaciones diferenciales}

Verifique si las siguientes expresiones son soluciones de las EDO dadas.

\begin{enumerate}[label=\alph*.-]

    \item $y=Cxe^x$, con $C\in\mathbb{R}$, es solución de
    $y''-2y'+y=0$.

    \solucion

    Derivamos aplicando la regla del producto:
    \[
    y'=C\left(e^x+xe^x\right)=Ce^x(x+1),
    \]
    \[
    y''=C\left[e^x(x+1)+e^x\right]=Ce^x(x+2).
    \]

    Reemplazando en el primer miembro de la ecuación:
    \[
    y''-2y'+y
    =Ce^x(x+2)-2Ce^x(x+1)+Cxe^x,
    \]
    \[
    y''-2y'+y=Ce^x\left[(x+2)-2(x+1)+x\right]
    =Ce^x\cdot0=0.
    \]

    \[
    \boxed{
    y=Cxe^x\ \text{ es solución de } y''-2y'+y=0\ \text{ para todo } C\in\mathbb{R}
    }
    \]


    \item $y=e^{3x}\cos(2x)$ es solución de
    $\dfrac{d^2y}{dx^2}-6\dfrac{dy}{dx}+13y=0$.

    \solucion

    Primera derivada (regla del producto):
    \[
    \frac{dy}{dx}=3e^{3x}\cos(2x)-2e^{3x}\sen(2x)
    =e^{3x}\left[3\cos(2x)-2\sen(2x)\right].
    \]

    Segunda derivada:
    \[
    \frac{d^2y}{dx^2}
    =3e^{3x}\left[3\cos(2x)-2\sen(2x)\right]
    +e^{3x}\left[-6\sen(2x)-4\cos(2x)\right],
    \]
    \[
    \frac{d^2y}{dx^2}=e^{3x}\left[5\cos(2x)-12\sen(2x)\right].
    \]

    Reemplazando:
    \[
    \frac{d^2y}{dx^2}-6\frac{dy}{dx}+13y
    =e^{3x}\Big[\left(5\cos2x-12\sen2x\right)
    -\left(18\cos2x-12\sen2x\right)
    +13\cos2x\Big].
    \]

    Agrupando términos:
    \[
    \cos(2x):\ 5-18+13=0,
    \qquad
    \sen(2x):\ -12+12=0.
    \]

    \[
    \boxed{
    y=e^{3x}\cos(2x)\ \text{ es solución de }
    \frac{d^2y}{dx^2}-6\frac{dy}{dx}+13y=0
    }
    \]


    \item $x^2+y^2=2C$ es solución de $\dfrac{dy}{dx}=-\dfrac xy$, donde
    $C\in\mathbb{R}$.

    \solucion

    La expresión define a $y$ implícitamente como función de $x$.
    Derivamos implícitamente respecto de $x$:
    \[
    2x+2y\frac{dy}{dx}=0
    \quad\Rightarrow\quad
    \frac{dy}{dx}=-\frac{x}{y}\qquad(y\neq0).
    \]

    Se obtiene exactamente la ecuación diferencial dada.

    \[
    \boxed{
    x^2+y^2=2C\ \text{ es solución (implícita) de } \frac{dy}{dx}=-\frac xy
    }
    \]

    Observación: para que la expresión represente una curva real se
    necesita $C>0$; en ese caso se trata de la familia de circunferencias
    centradas en el origen de radio $\sqrt{2C}$.

\end{enumerate}

\volverindice
\newpage

% ======================== EJERCICIO 17 ========================

\section{Ecuaciones diferenciales de variables separables}

Resuelve las siguientes ecuaciones diferenciales de variables
separables.

\begin{enumerate}[label=\alph*.-]

    \item
    \[
    \frac{dy}{dx}=\frac xy
    \]

    \solucion

    Separamos las variables e integramos:
    \[
    y\,dy=x\,dx
    \quad\Rightarrow\quad
    \int y\,dy=\int x\,dx
    \quad\Rightarrow\quad
    \frac{y^2}{2}=\frac{x^2}{2}+C.
    \]

    Multiplicando por 2 y llamando $K=2C$:
    \[
    \boxed{
    y^2-x^2=K,\qquad K\in\mathbb{R}
    }
    \]

    La solución general es una familia de hipérbolas equiláteras (y, para
    $K=0$, el par de rectas $y=\pm x$).


    \item
    \[
    y'=\frac{\ln x}{xy},\qquad y(1)=2
    \]

    \solucion

    Separamos las variables:
    \[
    y\,dy=\frac{\ln x}{x}\,dx.
    \]

    Integrando (en el segundo miembro, con $u=\ln x$, $du=\frac1x\,dx$):
    \[
    \frac{y^2}{2}=\frac{(\ln x)^2}{2}+C.
    \]

    Aplicamos la condición inicial $y(1)=2$:
    \[
    \frac{4}{2}=\frac{(\ln1)^2}{2}+C
    \quad\Rightarrow\quad
    C=2.
    \]

    Entonces $y^2=(\ln x)^2+4$. Como $y(1)=2>0$, tomamos la raíz positiva:
    \[
    \boxed{
    y=\sqrt{(\ln x)^2+4},\qquad x>0
    }
    \]


    \item
    \[
    \frac{dy}{dx}-ky=0,\qquad k\in\mathbb{R}\ \wedge\ k\neq0
    \]

    \solucion

    Reescribimos la ecuación como $\dfrac{dy}{dx}=ky$ y separamos las
    variables (para $y\neq0$):
    \[
    \frac{dy}{y}=k\,dx
    \quad\Rightarrow\quad
    \int\frac{dy}{y}=\int k\,dx
    \quad\Rightarrow\quad
    \ln|y|=kx+C.
    \]

    Despejando:
    \[
    |y|=e^{kx+C}=e^C e^{kx}
    \quad\Rightarrow\quad
    y=\pm e^Ce^{kx}.
    \]

    Llamando $A=\pm e^C$, e incluyendo la solución $y=0$ (que también
    verifica la ecuación) con $A=0$:
    \[
    \boxed{
    y=Ae^{kx},\qquad A\in\mathbb{R}
    }
    \]

    Es el modelo de crecimiento ($k>0$) o decaimiento ($k<0$) exponencial.


    \item
    \[
    y'-y=0,\qquad y(0)=2
    \]

    \solucion

    Es el caso particular $k=1$ del inciso anterior:
    \[
    \frac{dy}{y}=dx
    \quad\Rightarrow\quad
    \ln|y|=x+C
    \quad\Rightarrow\quad
    y=Ae^{x}.
    \]

    Aplicamos la condición inicial:
    \[
    y(0)=Ae^0=A=2.
    \]

    \[
    \boxed{
    y=2e^{x}
    }
    \]


    \item (*)
    \[
    x(x+2)\frac{dy}{dx}-y(2x+6)=0
    \]

    \solucion

    Separamos las variables:
    \[
    \frac{dy}{y}=\frac{2x+6}{x(x+2)}\,dx.
    \]

    Descomponemos en fracciones parciales:
    \[
    \frac{2x+6}{x(x+2)}=\frac Ax+\frac{B}{x+2}
    \quad\Rightarrow\quad
    2x+6=A(x+2)+Bx.
    \]

    Con $x=0$: $6=2A\Rightarrow A=3$.
    Con $x=-2$: $2=-2B\Rightarrow B=-1$.

    Integrando:
    \[
    \ln|y|=3\ln|x|-\ln|x+2|+C=\ln\left|\frac{x^3}{x+2}\right|+C.
    \]

    \[
    \boxed{
    y=\frac{Kx^3}{x+2},\qquad K\in\mathbb{R}
    }
    \]


    \item (*)
    \[
    xy-(1+x^2)y'=0,\qquad y(0)=1
    \]

    \solucion

    Reescribimos como $(1+x^2)\dfrac{dy}{dx}=xy$ y separamos:
    \[
    \frac{dy}{y}=\frac{x}{1+x^2}\,dx.
    \]

    Integrando (en el segundo miembro, con $u=1+x^2$, $du=2x\,dx$):
    \[
    \ln|y|=\frac12\ln\left(1+x^2\right)+C
    \quad\Rightarrow\quad
    y=K\sqrt{1+x^2}.
    \]

    Aplicamos la condición inicial:
    \[
    y(0)=K\sqrt1=K=1.
    \]

    \[
    \boxed{
    y=\sqrt{1+x^2}
    }
    \]

\end{enumerate}

\volverindice
\newpage

\end{document}
